% ============================================================================
% Chapter 30 --- Coexistence patterns library
% ----------------------------------------------------------------------------
% Purpose : the recurring dual-stack situations encountered during a
%           migration, with their bridging patterns. Closes Part III.
% Page target: ~5 pages.
% Dependencies: Ch. 18-27 (each wave invokes coexistence as a principle).
% ============================================================================

\chapter{Coexistence patterns library}
\label{ch:coexistence-patterns}

The wave chapters of Part~III each declare a coexistence pattern as
the operational state of the wave. Several of those patterns recur
across waves and across institutions; this chapter consolidates the
six most frequently encountered, presents their bridging mechanisms,
and notes the anti-patterns that recur with them. The patterns are
described at the level of operational principle rather than at the
level of specific tooling, since their realisation varies with the
products the institution operates.

%-----------------------------------------------------------------------------
\section{Pattern: incumbent enterprise directory \texorpdfstring{$\leftrightarrow$}{<->} institutional IdP}
\label{sec:cp-ad-keycloak}

The pattern in which the incumbent enterprise directory remains
the authoritative source for staff identities while a new
institutional identity provider (an open-source IdP or comparable)
handles federation, role activation, and downstream OIDC
integrations.

\paragraph{Bridging mechanism.} The incumbent enterprise
directory, supplied by the institution's incumbent identity-suite
vendor, continues to authenticate users at the institutional
perimeter. The new IdP federates against the directory through
standard SAML2 or OIDC bridges. Attributes from the directory ---
group membership, role, organisational unit --- are mapped to the
institutional schema (D7.7) at the boundary of the new IdP rather
than propagated into downstream applications in their native form.

\paragraph{Synchronisation direction.} The pattern is
unidirectional from the incumbent directory toward the new IdP for
the duration of Wave~1. Reverse synchronisation (changes made in
the new IdP flowing back to the directory) is generally avoided to
keep the directory's authoritative status unambiguous.

\paragraph{Drift detection and exit window.} Discrepancies
between the directory and the new IdP are surfaced through the
observability layer. The pattern's exit window is the period
during which the institution can return to directory-only
operation if the new IdP fails; the window is typically aligned
with the legacy IdP retirement schedule.

%-----------------------------------------------------------------------------
\section{Pattern: incumbent endpoint-management suite \texorpdfstring{$\leftrightarrow$}{<->} declarative configuration}
\label{sec:cp-sccm-ansible}

The pattern in which the institution's existing endpoint
configuration management (the incumbent vendor's unified endpoint
management suite) coexists with declarative configuration
(a declarative configuration-management toolchain) for
the populations the new approach covers.

\paragraph{Bridging mechanism.} Each population is governed by a
single source. Staff desktop endpoints typically remain on the
incumbent-suite side; servers and research workstations typically
move to the declarative-configuration side; other endpoint classes
can sit on either side depending on the institution's mature
tooling. The boundary between the populations is documented and
stable for the duration of the coexistence.

\paragraph{Configuration overlap.} Where a single endpoint must be
visible to both systems (rare, but it occurs in shared-use scenarios),
the institution declares which system is authoritative for which
configuration domain --- typically the incumbent suite for
compliance policy and the declarative toolchain for application
configuration.

\paragraph{Posture export.} The unified posture schema of D11.6
ingests posture from both sides. Policies in
\trchap{ch:policy}{} are authored against the unified schema and
do not depend on which configuration management produced the
signal.

%-----------------------------------------------------------------------------
\section{Pattern: per-volume-priced commercial SIEM \texorpdfstring{$\leftrightarrow$}{<->} open-source store}
\label{sec:cp-splunk-opensearch}

The pattern in which a per-volume-priced commercial SIEM operates
alongside an open-source store, with the institution managing the
boundary deliberately for operational and financial reasons.

\paragraph{Bridging mechanism.} The collector layer of Wave~2
emits the canonical TAC schema to both stores. The commercial SIEM
ingests the events that the SOC actively investigates within the
investigation window the institution chooses to fund; the
open-source store ingests the longer-tail retention that the
commercial SIEM's per-GB pricing makes unattractive. Detection
content is authored in both stores where the institution's risk
profile warrants duplication.

\paragraph{Per-source routing.} Sources whose events are most
operationally consumed by the SOC (identity, network attachment,
endpoint posture) are routed to the commercial SIEM in full;
sources whose events serve audit-replay rather than routine
investigation (verbose flow records, low-severity informational
events) are routed primarily to the open-source store with a
sampled stream to the commercial SIEM for context.

\paragraph{Migration of dashboards.} Commercial-SIEM dashboards are
migrated incrementally to the open-source store and its
visualisation layer when the institution chooses to retire a
commercial-SIEM content element. The choice to retire is informed
by financial review rather than by technical inability.

%-----------------------------------------------------------------------------
\section{Pattern: incumbent commercial hypervisor \texorpdfstring{$\leftrightarrow$}{<->} open virtualisation substrate}
\label{sec:cp-vmware-k8s}

The pattern in which the incumbent commercial hypervisor platform
continues to host the institution's existing workloads while new
workloads are provisioned on an open virtualisation substrate
(a Kubernetes-native VM platform or an open-source IaaS), with
explicit governance of which platform hosts which workload.

\paragraph{Bridging mechanism.} The governance layer (Wave~8
classification) decides per workload class which platform is
appropriate: existing legacy workloads stay on the commercial
hypervisor through their natural lifecycle; new workloads provision
on the open substrate by default; specific workload classes that
benefit from the commercial platform's operational features (live
migration, mature DR) remain on it by explicit decision.

\paragraph{Live-migration boundary.} Live migration between
substrates is technically possible (the open substrate can import
the commercial platform's VMs, and supports VM-level migration) but
is treated as a project rather than a routine operation. Cold
migration on workload retirement / refresh is the more common
pattern.

\paragraph{Skills coexistence.} Operational teams need both
skill sets during the coexistence, and the institution's staffing
profile must carry both for the duration. The skill investment is
typically an explicit programme item rather than a passive
expectation.

%-----------------------------------------------------------------------------
\section{Pattern: enterprise storage array \texorpdfstring{$\leftrightarrow$}{<->} open-source distributed storage}
\label{sec:cp-netapp-ceph}

The pattern in which an enterprise storage array from a commercial
vendor continues to host the institution's existing data corpora
while an open-source distributed storage cluster provides storage
for new workloads and for the research-scale capacity the
enterprise platform is not cost-effective for.

\paragraph{Bridging mechanism.} The classification scheme of
D10.4 determines which platform hosts which data class. Enterprise
storage typically retains administrative and compliance data; the
open-source cluster typically holds research data, container
persistent volumes, and the cold-tier object storage that the
observability layer (D12.3) needs.

\paragraph{Migration of pools.} Migration from one platform to
the other occurs at workload boundaries: a workload migrating
between platforms (Wave~8) typically migrates its storage at the
same time. Storage-only migration without workload migration is
operationally costly and is reserved for specific drivers
(end-of-support, cost reset, disclosure-regime change).

\paragraph{Data-integrity validation.} A migration between
platforms validates data integrity through documented checksums
or content-addressable verification. The validation is part of
the conformance suite of \trchap{ch:conformance}{}
(data-portability category).

%-----------------------------------------------------------------------------
\section{Pattern: eduroam \texorpdfstring{$\leftrightarrow$}{<->} TAC-aware Wi-Fi}
\label{sec:cp-eduroam-tac}

Where the institution participates in a federated academic-roaming
service (in the research-and-education community, eduroam), this is
the pattern in which that participation continues for
inter-institutional roaming while a TAC-aware RADIUS attachment
governs intra-institutional connections.

\paragraph{Bridging mechanism.} The institution's RADIUS service
recognises both authentication paths. eduroam-routed
authentications (visiting researchers from peer institutions,
staff roaming abroad) follow the eduroam attribute schema
unchanged. Local authentications (staff and students at the
institution's sites) carry the institutional schema with TAC-
relevant attributes that the policy plane consumes.

\paragraph{Attribute schema preservation.} The federated-roaming
technical profile (eduroam, where used) is treated as inviolate:
extending it with institution-specific attributes for
inter-institutional sessions would break federation. The
institutional attributes are added only on the non-roaming
authentication path.

\paragraph{Roaming policy.} Visiting researchers receive the
default eduroam treatment (Internet access from a guest zone)
unless an explicit institutional arrangement (a hosted research
project, a long-stay visitor) provisions a TAC for them.

%-----------------------------------------------------------------------------
\section{Common anti-patterns}
\label{sec:cp-anti-patterns}

The patterns above describe how coexistence is managed. The
anti-patterns below describe how it goes wrong.

\paragraph{Long-running dual stack without exit date.} A coexistence
period whose end is not declared at the wave's start tends to
become permanent. The institution operates two of everything
indefinitely, and the cost grows. The mitigation is to declare the
exit date at wave open, to track it through the programme
governance, and to escalate when the date drifts.

\paragraph{Unspecified bridging ownership.} The bridge between two
systems requires operational ownership. When ownership is implicit,
the bridge fails silently --- a synchronisation job stops, a
metadata refresh expires, a policy mapping drifts --- and the
discovery is during an incident. The mitigation is to name the
bridge owner explicitly at wave open, with documented escalation
path.

\paragraph{Bidirectional sync without conflict resolution.} A
bridge that synchronises in both directions without a documented
conflict-resolution rule produces silent data loss when the two
sides diverge. The mitigation is unidirectional sync wherever
possible; where bidirectional is unavoidable, an explicit
conflict-resolution rule documented and tested before production
deployment.

\paragraph{Coexistence as a permanent feature.} Some institutions
stabilise on a coexistence pattern and treat it as the architecture
rather than as a transition. The pattern is not always wrong (a
hybrid open-source federation-and-IdP identity arrangement is a
legitimate end-state, for example), but it should be a deliberate
decision
rather than an accident. The mitigation is the periodic review
under the sovereignty grid: a coexistence that has not been
revisited in twenty-four months is reviewed against its current
costs and benefits.

\medskip

\noindent The patterns and anti-patterns above are the principal
operational shapes of coexistence in the playbook. Their explicit
documentation in the institution's wave plans turns coexistence
from an unplanned consequence of migration into a deliberate
engineering artefact.
